% ===========================
% capitulo_implementacion_logico.tex
% ===========================

\section{Módulo Lógico}% chktex 24
\label{cap:implementacion-modulo-logico}

\subsection{Introducción}
Este capítulo documenta la implementación del \emph{Módulo Lógico}, responsable de orquestar la optimización evolutiva de las masas y la evaluación de estabilidad dinámica del sistema de dos/tres cuerpos. El diseño actual concentra la coordinación de flujos, el motor genético en modo \emph{streaming}, la evaluación de fitness basada en Lyapunov y la lógica auxiliar de \emph{re-seeding}, todo ello instrumentado con \texttt{time\_block}/\texttt{time\_section} para alimentar la telemetría global.

\subsection{Arquitectura de Alto Nivel}
El módulo se compone de cuatro bloques principales:
\begin{enumerate}
    \item \textbf{Controlador continuo (\texttt{controller.py})}: coordina las iteraciones del GA, respeta presupuestos de tiempo/evaluaciones, consulta horizontes corto/largo y registra métricas por época.
    \item \textbf{Motor genético (\texttt{ga.py})}: implementa un GA incremental basado en \texttt{pymoo} con degradación a un sampler interno; soporta distribuciones Beta para masas y reseeds locales.
    \item \textbf{Evaluador de fitness (\texttt{fitness.py})}: evalúa candidatos con el estimador de Lyapunov, realiza caching jerárquico y aplica penalizaciones de cuasi-periodicidad.
    \item \textbf{Modificador de parámetros (\texttt{parameters.py})}: encapsula reglas de re-seeding durante el estancamiento. El archivo \texttt{\_\_init\_\_.py} re-exporta la API pública.
\end{enumerate}

\subsection{Stack tecnológico}
\begin{itemize}
    \item \textbf{Python 3.10+} como lenguaje principal, con tipado estático opcional.
    \item \textbf{NumPy} para operaciones numéricas y manejo vectorial en el GA y el evaluador de fitness.
    \item \textbf{pymoo} (opcional) para algoritmos evolutivos; el código se degrada con un sampler propio si la librería no está disponible.
    \item \textbf{Infraestructura interna}: \texttt{Config}, \texttt{HierarchicalCache}, \texttt{MetricsLogger}, \texttt{perf\_timings.timers} y el módulo de simulación (\texttt{ReboundSim}, \texttt{LyapunovEstimator}).
\end{itemize}

% =========================================
% DOCUMENTACIóN POR SCRIPT
% =========================================

% ---------- controller.py ----------
\subsection{\texttt{controller.py}\-- Controlador de Optimización Continua}
\paragraph{Descripción general}
\texttt{ContinuousOptimizationController} es la puerta de entrada del módulo. Instancia \texttt{StreamingGA}, \texttt{FitnessEvaluator}, \texttt{HierarchicalCache} y \texttt{MetricsLogger}; coordina las evaluaciones a horizonte corto/largo, aplica penalizaciones por falta de periodicidad y captura telemetría detallada por época.

\subsubsection{\texttt{class ContinuousOptimizationController}}
\paragraph{Propósito} Coordinar el ciclo de optimización evolutiva, generando candidatos, evaluando su fitness, reseedeando la población cuando hay estancamiento y devolviendo el mejor resultado junto con un resumen métrico completo.

\subsubsection*{\texttt{\_\_init\_\_(cfg:\ Config, logger:\ Optional[Any] = None) -> None}}

\begin{table}[H]
  \centering
  \caption{Parómetros de \texttt{ContinuousOptimizationController.\_\_init\_\_}}
  \small
  \begin{tabularx}{\textwidth}{@{}p{0.27\textwidth}p{0.33\textwidth}X@{}}
    \toprule
    \textbf{Parómetro} & \textbf{Tipo} & \textbf{Descripción} \\
    \midrule
    \texttt{cfg} & \texttt{Config} & Configuración global (presupuestos, radios, masas iniciales, etc.). \\
    \texttt{logger} & \texttt{Any | None} & Registrador opcional; si es \texttt{None}, se usa \texttt{logging.getLogger}. \\
    \bottomrule
  \end{tabularx}
\end{table}

\begin{table}[H]
  \centering
  \caption{Valor de retorno de \texttt{ContinuousOptimizationController.\_\_init\_\_}}
  \small
  \begin{tabularx}{\textwidth}{@{}p{0.35\textwidth}X@{}}
    \toprule
    \textbf{Tipo} & \textbf{Descripción} \\
    \midrule
    \texttt{None} & Constructor; inicializa métricas y modificadores auxiliares. \\
    \bottomrule
  \end{tabularx}
\end{table}

\begin{listing}[H]
\caption{Bloque de código de \texttt{ContinuousOptimizationController.\_\_init\_\_}}
\begin{adjustbox}{max width=\textwidth}
\begin{minipage}{\textwidth}
\begin{minted}[fontsize=\small, breaklines, breakanywhere]{python}
class ContinuousOptimizationController:
    def __init__(self, cfg: Config, logger: Optional[Any] = None) -> None:
        import logging

        self.cfg = cfg
        self.logger = logger or logging.getLogger("master_two_body_opt")
        self.metrics = MetricsLogger()
        self._modifier = ParameterModifier(cfg)
\end{minted}
\end{minipage}
\end{adjustbox}
\end{listing}

\subsubsection*{\texttt{run~() -> Dict[str, Any]}}
\textbf{Propósito.} Ejecutar la optimización completa respetando presupuestos (\texttt{time\_budget\_s}, \texttt{eval\_budget}, \texttt{max\_epochs}), intercalando evaluaciones a horizonte corto/largo, registrando métricas por época y devolviendo el mejor candidato encontrado con sus estadísticas.

\begin{table}[H]
  \centering
  \caption{Valor de retorno de \texttt{ContinuousOptimizationController.run}}
  \small
  \begin{tabularx}{\textwidth}{@{}p{0.35\textwidth}X@{}}
    \toprule
    \textbf{Tipo} & \textbf{Descripción} \\
    \midrule
    \texttt{Dict[str, Any]} & \{\texttt{``status''}, \texttt{``best''}, \texttt{``evals''}, \texttt{``epochs''}\}; \texttt{best} incluye masas, \texttt{lambda} y \texttt{fitness}. \\
    \bottomrule
  \end{tabularx}
\end{table}

\begin{listing}[H]
\caption{Fragmento principal de \texttt{ContinuousOptimizationController.run}}
\begin{adjustbox}{max width=0.85\textwidth}
\begin{minipage}{\linewidth}
{\setstretch{0.95}
\begin{adjustbox}{max width=\textwidth}
\begin{minipage}{\textwidth}
\begin{minted}[fontsize=\scriptsize, breaklines, breakanywhere]{python}
    def run(self) -> Dict[str, Any]:
        start = time.time()
        cache = HierarchicalCache(
            approx_max=self.cfg.cache_approx_max, exact_max=self.cfg.cache_exact_max
        )
        evaluator = FitnessEvaluator(cache, self.cfg, logger=self.logger)
        ga = StreamingGA(self.cfg, logger=self.logger)

        best: Optional[Tuple[float, ...]] = None
        best_fit = -float("inf")
        best_lambda_value: Optional[float] = None
        stagnation_counter = 0
        radius = self.cfg.local_radius
        epoch = 0
        run_start = time.time()
        if self.logger:
            self.logger.info(
                "Starting optimization | pop=%d | dims=%d | time_budget=%.1fs | eval_budget=%d",
                self.cfg.pop_size,
                len(getattr(self.cfg, "mass_bounds", tuple())),
                self.cfg.time_budget_s,
                self.cfg.eval_budget,
            )

        while True:
            epoch_start = time.time()
            now = time.time()
            if now - start > self.cfg.time_budget_s or self.metrics.eval_count >= self.cfg.eval_budget:
                break
            if epoch >= self.cfg.max_epochs:
                break

            X = ga.current_population()
            fits, short_details = evaluator.evaluate_batch(
                X, horizon="short", epoch=epoch, return_details=True
            )
            self.metrics.eval_count += len(X)
            short_evals = len(X)
            long_evals = 0
            epoch_best_lambda_short: Optional[float] = None
            epoch_best_fitness_short: Optional[float] = None
            epoch_best_lambda_global: Optional[float] = best_lambda_value
            epoch_best_fitness_global: Optional[float] = (
                float(best_fit) if best is not None and best_fit != -float("inf") else None
            )
            epoch_source = None
            epoch_candidate_short: Optional[Tuple[float, ...]] = None

            if len(fits) > 0:
                idx = int(max(range(len(fits)), key=lambda i: fits[i]))
                epoch_best_fit = float(fits[idx])
                epoch_best = X[idx]
                detail_short: Dict[str, Any] = short_details[idx] if idx < len(short_details) else {}
                lambda_short_val = detail_short.get("lambda") if isinstance(detail_short, dict) else None
                if lambda_short_val is None:
                    lambda_short_val = -epoch_best_fit
                penalty_short_val = (
                    detail_short.get("penalty") if isinstance(detail_short, dict) else None
                )
                epoch_best_lambda_short = float(lambda_short_val)
                epoch_best_fitness_short = epoch_best_fit
                epoch_candidate_short = tuple(epoch_best)
                if epoch_best_fit > best_fit:
                    best_fit = epoch_best_fit
                    best = epoch_best
                    best_lambda_value = float(lambda_short_val)
                    stagnation_counter = 0
                    epoch_source = "short"
                    epoch_best_lambda_global = best_lambda_value
                    epoch_best_fitness_global = best_fit
                    if self.logger:
                        penalty_suffix = (
                            f" | penalty={penalty_short_val:.6f}"
                            if penalty_short_val is not None
                            else ""
                        )
                        self.logger.info(
                            "Epoch %d | new global best (short) lambda=%.6f | fitness=%.6f%s | masses=%s",
                            epoch,
                            best_lambda_value,
                            best_fit,
                            penalty_suffix,
                            tuple(round(float(v), 6) for v in epoch_best),
                        )
                else:
                    stagnation_counter += 1

            try:
                import numpy as _np

                ga._last_F = _np.array(fits, dtype=float).reshape(-1, 1)
            except Exception:
                ga._last_F = None
            ga.step(generations=max(1, self.cfg.n_gen_step), epoch=epoch)

            if len(fits) > 0 and self.cfg.top_k_long > 0:
                k = min(self.cfg.top_k_long, len(fits))
                idxs = sorted(range(len(fits)), key=lambda i: fits[i], reverse=True)[:k]
                topX = [X[i] for i in idxs]
                long_fits, long_details = evaluator.evaluate_batch(
                    topX, horizon="long", epoch=epoch, return_details=True
                )
                self.metrics.eval_count += len(topX)
                long_evals = len(topX)
                for xi, f_long, detail_long in zip(topX, long_fits, long_details):
                    if f_long > best_fit:
                        lambda_long_val = (
                            detail_long.get("lambda") if isinstance(detail_long, dict) else None
                        )
                        if lambda_long_val is None:
                            lambda_long_val = -float(f_long)
                        penalty_long_val = (
                            detail_long.get("penalty") if isinstance(detail_long, dict) else None
                        )
                        best_fit = float(f_long)
                        best = xi
                        best_lambda_value = float(lambda_long_val)
                        epoch_source = "long"
                        stagnation_counter = 0
                        epoch_best_lambda_global = best_lambda_value
                        epoch_best_fitness_global = best_fit
                        if self.logger:
                            penalty_suffix = (
                                f" | penalty={penalty_long_val:.6f}"
                                if penalty_long_val is not None
                                else ""
                            )
                            self.logger.info(
                                "Epoch %d | new global best (long) lambda=%.6f | fitness=%.6f%s | masses=%s",
                                epoch,
                                best_lambda_value,
                                best_fit,
                                penalty_suffix,
                                tuple(round(float(v), 6) for v in xi),
                            )

            if stagnation_counter >= self.cfg.stagnation_window:
                if self.logger:
                    self.logger.info("Stagnation detected; reseeding around best candidate.")
                radius = self._modifier.on_stagnation(ga, best, radius)
                self.metrics.reseeds += 1
                stagnation_counter = 0

            self.metrics.record_epoch(
                {
                    "epoch": epoch,
                    "best_lambda_short": epoch_best_lambda_short,
                    "best_fitness_short": epoch_best_fitness_short,
                    "best_candidate_short": list(epoch_candidate_short) if epoch_candidate_short else None,
                    "best_lambda_global": epoch_best_lambda_global,
                    "best_fitness_global": epoch_best_fitness_global,
                    "best_candidate_global": list(best) if best else None,
                    "evaluations": {
                        "short": short_evals,
                        "long": long_evals,
                        "total": self.metrics.eval_count,
                    },
                    "stagnation_counter": stagnation_counter,
                    "radius": radius,
                    "epoch_time_s": time.time() - epoch_start,
                    "source": epoch_source,
                    "timestamp": time.time(),
                }
            )

            epoch += 1
            self.metrics.epochs = epoch

        self.metrics.mark_phase("completed")
        if best is not None:
            final_lambda = (
                float(best_lambda_value) if best_lambda_value is not None else -float(best_fit)
            )
        else:
            final_lambda = None
        if self.logger:
            self.logger.info(
                "Optimization completed | epochs=%d | evals=%d | best lambda=%s | wall=%.1fs",
                epoch,
                self.metrics.eval_count,
                f"{final_lambda:.6f}" if final_lambda is not None else "N/A",
                time.time() - run_start,
            )

        best_payload: Dict[str, Any] = {
            "masses": list(best) if best else None,
            "lambda": final_lambda,
            "fitness": float(best_fit) if best is not None else None,
        }
        if best:
            if len(best) > 0:
                best_payload["m1"] = float(best[0])
            if len(best) > 1:
                best_payload["m2"] = float(best[1])
            if len(best) > 2:
                best_payload["m3"] = float(best[2])

        return {
            "status": "completed",
            "best": best_payload,
            "evals": self.metrics.eval_count,
            "epochs": epoch,
        }
\end{minted}
\end{minipage}
\end{adjustbox}


}
\end{minipage}
\end{adjustbox}
\end{listing}

Durante cada época el controlador actualiza el \texttt{MetricsLogger} con lambdas corto/largo, penalizaciones, conteo de reseeds y tiempos, habilitando la generación de reportes de convergencia sin re-ejecutar simulaciones.



\begin{listing}[H]
\caption{Ejemplo de uso de \texttt{ContinuousOptimizationController}}
\begin{adjustbox}{max width=\textwidth}
\begin{minipage}{\textwidth}
\begin{minted}[fontsize=\small, breaklines, breakanywhere]{python}
from two_body.core.config import Config
from two_body.core.telemetry import setup_logger
from two_body.logic.controller import ContinuousOptimizationController

cfg = Config(max_epochs=1, time_budget_s=1.0, eval_budget=10)
logger = setup_logger()
controller = ContinuousOptimizationController(cfg, logger=logger)
result = controller.run()
print("Resultado parcial:", result)
\end{minted}
\end{minipage}
\end{adjustbox}
\end{listing}

% ---------- ga.py ----------
\subsection{\texttt{ga.py} \- Motor Genótico en Streaming}
\paragraph{Descripción general}
\texttt{StreamingGA} envuelve la configuración del algoritmo genético de \texttt{pymoo} (\texttt{GA}) para el caso unidimensional de minimización (mLCE negativo). Expone una interfaz incremental (\texttt{ask}/\texttt{tell}) y se degrada a un sampler aleatorio con recombinación propia si \texttt{pymoo} no estó instalado o falla el \emph{setup}.

\subsubsection{\texttt{class StreamingGA}}
\paragraph{Propósito} Generar poblaciones candidato, recibir evaluaciones y evolucionar la exploración según los hiperparámetros definidos en \texttt{Config} (población, operadores, distribución de masas, elitismo), manteniendo soporte para reseeds locales y degradación controlada.

\subsubsection*{\texttt{\_\_init\_\_(cfg:\ Config, logger:\ Optional[Any] = None) -> None}}

\begin{table}[H]
  \centering
  \caption{Parómetros de \texttt{StreamingGA.\_\_init\_\_}}
  \small
  \begin{tabularx}{\textwidth}{@{}p{0.27\textwidth}p{0.33\textwidth}X@{}}
    \toprule
    \textbf{Parómetro} & \textbf{Tipo} & \textbf{Descripción} \\
    \midrule
    \texttt{cfg} & \texttt{Config} & Proporciona límites de masas, semillas, tasas de cruce/mutación y tamaño poblacional. \\
    \texttt{logger} & \texttt{Any | None} & Registrador opcional para mensajes informativos o de degradación. \\
    \bottomrule
  \end{tabularx}
\end{table}

\begin{listing}[H]
\caption{Bloque de código de \texttt{StreamingGA.\_\_init\_\_}}
\begin{adjustbox}{max width=\textwidth}
\begin{minipage}{0.85\textwidth}
\begin{minted}[fontsize=\scriptsize, breaklines, breakanywhere]{python}
class StreamingGA:
    def __init__(self, cfg: Config, logger: Optional[Any] = None) -> None:
        import logging

        self.cfg = cfg
        self.logger = logger or logging.getLogger("master_two_body_opt")
        self._asked_pop = None
        self._asked_X: Optional[Any] = None
        self._last_F: Optional[Any] = None
        self._rng = random.Random(cfg.seed)
        self._reseed_center: Optional[Tuple[float, ...]] = None
        self._reseed_radius: float = cfg.local_radius
        self._bounds = [tuple(map(float, b)) for b in cfg.mass_bounds]
        self._dim = len(self._bounds)
        self._mass_distribution = getattr(cfg, "mass_distribution", "uniform").lower()
        if self._mass_distribution not in {"uniform", "beta"}:
            if self.logger:
                self.logger.warning(
                    "Mass distribution '%s' no soportada; se usara uniforme.",
                    self._mass_distribution,
                )
            self._mass_distribution = "uniform"
        self._beta_alpha = self._normalize_beta_param(
            getattr(cfg, "mass_beta_alpha", 1.0),
            "mass_beta_alpha",
        )
        self._beta_beta = self._normalize_beta_param(
            getattr(cfg, "mass_beta_beta", 1.0),
            "mass_beta_beta",
        )

        try:
            from pymoo.algorithms.soo.nonconvex.ga import GA
            from pymoo.core.problem import Problem
            from pymoo.core.sampling import Sampling
            from pymoo.operators.crossover.sbx import SBX
            from pymoo.operators.mutation.pm import PolynomialMutation
            from pymoo.operators.selection.tournament import TournamentSelection

            class _TimedSBX(SBX):
                def do(self, *args, **kwargs):
                    with time_block("crossover", extra={"impl": "SBX"}):
                        return super().do(*args, **kwargs)

            class _TimedPolynomialMutation(PolynomialMutation):
                def do(self, *args, **kwargs):
                    with time_block("mutation", extra={"impl": "PolynomialMutation"}):
                        return super().do(*args, **kwargs)

            class _TimedTournamentSelection(TournamentSelection):
                def do(self, *args, **kwargs):
                    with time_block("selection_tournament", extra={"impl": "pymoo"}):
                        return super().do(*args, **kwargs)

            class _MassProblem(Problem):
                def __init__(self, lows: Sequence[float], highs: Sequence[float]):
                    import numpy as _np

                    super().__init__(
                        n_var=len(lows),
                        n_obj=1,
                        n_constr=0,
                        xl=_np.array(lows, dtype=float),
                        xu=_np.array(highs, dtype=float),
                    )

                def _evaluate(self, X, out, *args, **kwargs):
                    import numpy as _np

                    out["F"] = _np.zeros((len(X), 1), dtype=float)

            class _BetaSampling(Sampling):
                def __init__(
                    self,
                    lows: Sequence[float],
                    highs: Sequence[float],
                    alpha: Sequence[float],
                    beta: Sequence[float],
                    rng: random.Random,
                ) -> None:
                    super().__init__()
                    self._lows = tuple(float(v) for v in lows)
                    self._highs = tuple(float(v) for v in highs)
                    self._alpha = tuple(float(v) for v in alpha)
                    self._beta = tuple(float(v) for v in beta)
                    self._rng = rng

                def _do(self, problem, n_samples, **kwargs):
                    import numpy as _np

                    n_dim = len(self._lows)
                    X = _np.empty((n_samples, n_dim), dtype=float)
                    for i in range(n_samples):
                        for j in range(n_dim):
                            alpha = self._alpha[j % len(self._alpha)]
                            beta = self._beta[j % len(self._beta)]
                            beta_val = self._rng.betavariate(alpha, beta)
                            X[i, j] = self._lows[j] + beta_val * (
                                self._highs[j] - self._lows[j]
                            )
                    return X

            lows = [b[0] for b in self._bounds]
            highs = [b[1] for b in self._bounds]
            self._problem = _MassProblem(lows, highs)
            crossover = _TimedSBX(prob=max(0.0, min(1.0, float(cfg.crossover))))
            mutation = _TimedPolynomialMutation(eta=cfg.mutation_eta, prob=cfg.mutation_prob)
            selection = _TimedTournamentSelection(func_comp=lambda a, b: a.F[0] >= b.F[0])
            sampling = _BetaSampling(lows, highs, self._beta_alpha, self._beta_beta, self._rng)
            ga_kwargs = {
                "pop_size": cfg.pop_size,
                "sampling": sampling if self._mass_distribution == "beta" else None,
                "crossover": crossover,
                "mutation": mutation,
                "selection": selection,
                "eliminate_duplicates": False,
            }
            self._algorithm = GA(**ga_kwargs)
            if cfg.selection != "tournament" and self.logger:
                self.logger.warning(
                    "Selection '%s' not supported explicitly; pymoo default will be used.",
                    cfg.selection,
                )
            self._algorithm.setup(self._problem, termination=None, seed=cfg.seed)
        except Exception as e:
            self._problem = None
            self._algorithm = None
            if self.logger:
                self.logger.warning(
                    "pymoo unavailable or setup failed: %s. Using internal sampler.", e
                )
            self._pop = self._init_random_population()
\end{minted}
\end{minipage}
\end{adjustbox}
\end{listing}

\subsubsection*{\texttt{current\_population~() -> List[Tuple[float, \dots]]}}
\textbf{Propósito.} Devolver la población actual. Con backend \texttt{pymoo} usa \texttt{ask()} y aplica reseeds locales; en modo degradado genera candidatos con \texttt{\_sample\_gene} (distribución uniforme/Beta) respetando los límites configurados.

\begin{table}[H]
  \centering
  \caption{Valor de retorno de \texttt{current\_population}}
  \small
  \begin{tabularx}{\textwidth}{@{}p{0.35\textwidth}X@{}}
    \toprule
    \textbf{Tipo} & \textbf{Descripción} \\
    \midrule
    \texttt{List[Tuple[float, \dots]]} & Secuencia de candidatos, cada uno con tantas masas como dimensiones configuradas. \\
    \bottomrule
  \end{tabularx}
\end{table}

\begin{listing}[H]
\caption{Bloque de código de \texttt{current\_population}}
\begin{adjustbox}{max width=0.65\textwidth}
\begin{minipage}{\textwidth}
\begin{minted}[fontsize=\small, breaklines, breakanywhere]{python}
    def current_population(self) -> List[Tuple[float, ...]]:
        import numpy as _np

        if self._algorithm is None:
            if getattr(self, "_pop", None) is None:
                bounds = self._bounds
                center = self._reseed_center
                radius = self._reseed_radius
                self._pop = [
                    tuple(
                        min(
                            max(
                                (
                                    center[idx]
                                    if center
                                    else self._sample_gene(idx, bounds)
                                )
                                + (self._rng.uniform(-radius, radius) if center else 0.0),
                                bounds[idx][0],
                            ),
                            bounds[idx][1],
                        )
                        for idx in range(self._dim)
                    )
                    for _ in range(self.cfg.pop_size)
                ]
            return list(self._pop)

        if self._asked_pop is not None and self._asked_X is not None:
            X = self._asked_X
        else:
            pop = self._algorithm.ask()
            X = pop.get("X")
            if self._reseed_center is not None and self._reseed_radius > 0:
                bounds = self._bounds
                center = self._reseed_center
                rad = self._reseed_radius
                for i in range(len(X)):
                    for j in range(self._dim):
                        X[i, j] = float(
                            _np.clip(
                                center[j] + self._rng.uniform(-rad, rad),
                                bounds[j][0],
                                bounds[j][1],
                            )
                        )
            pop.set("X", X)
            self._asked_pop = pop
            self._asked_X = X

        return [tuple(float(x[j]) for j in range(self._dim)) for x in X]
\end{minted}
\end{minipage}
\end{adjustbox}
\end{listing}

\subsubsection*{\texttt{step~(generations:\ float = 1.0) -> None}}
\vfill
\textbf{Propósito.} Avanzar el GA con las evaluaciones acumuladas en \texttt{\_last\_F}, instrumentando el backend (\texttt{pymoo} o degradado) y aplicando elitismo/reseed cuando corresponde.

\begin{listing}[H]
\caption{Fragmento de \texttt{StreamingGA.step}}
\begin{adjustbox}{max width=0.65\textwidth}
\begin{minipage}{\textwidth}
\begin{minted}[fontsize=\scriptsize, breaklines, breakanywhere]{python}
    @time_section(
        "ga_main",
        epoch=lambda self, generations=1.0, epoch=-1: epoch,
        extra=lambda self, generations=1.0, epoch=-1: {
            "generations": generations,
            "backend": "pymoo" if self._algorithm is not None else "fallback",
        },
    )
    def step(self, generations: float = 1.0, *, epoch: int = -1) -> None:
        if self._algorithm is None:
            if getattr(self, "_pop", None) is None:
                self._pop = self._init_random_population()
            bounds = self._bounds
            center = self._reseed_center
            radius = self._reseed_radius
            elites: List[Tuple[float, ...]] = []
            if (
                self.cfg.elitism > 0
                and self._last_F is not None
                and getattr(self, "_pop", None) is not None
            ):
                with time_block(
                    "selection_tournament",
                    extra={
                        "impl": "fallback",
                        "candidates": len(self._pop) if getattr(self, "_pop", None) else 0,
                    },
                ):
                    try:
                        scores = [float(val[0]) for val in list(self._last_F)]
                    except Exception:
                        scores = []
                    if scores:
                        ranked = sorted(
                            range(len(self._pop)),
                            key=lambda i: scores[i],
                            reverse=True,
                        )
                        for idx in ranked[: min(self.cfg.elitism, len(self._pop))]:
                            elites.append(self._pop[idx])
            new_pop: List[Tuple[float, ...]] = list(elites)
            while len(new_pop) < self.cfg.pop_size:
                if center is not None:
                    genes = [
                        min(
                            max(
                                center[idx] + self._rng.uniform(-radius, radius),
                                bounds[idx][0],
                            ),
                            bounds[idx][1],
                        )
                        for idx in range(self._dim)
                    ]
                else:
                    genes = self._sample_offspring(bounds)
                new_pop.append(tuple(genes))
            self._pop = new_pop
            self._asked_pop = None
            self._asked_X = None
            self._last_F = None
            return

        if self._asked_pop is None or self._asked_X is None:
            _ = self.current_population()
        if self._last_F is None:
            if self.logger:
                self.logger.debug("No fitness to tell; skipping GA step.")
            return
        try:
            self._asked_pop.set("F", self._last_F)
            self._algorithm.tell(infills=self._asked_pop)
            for _ in range(max(1, int(generations))):
                self._algorithm.next()
        finally:
            self._asked_pop = None
            self._asked_X = None
            self._last_F = None
\end{minted}
\end{minipage}
\end{adjustbox}
\end{listing}

\begin{listing}[H]
\caption{Ejemplo de inicialización de \texttt{StreamingGA}}
\begin{adjustbox}{max width=\textwidth}
\begin{minipage}{\textwidth}
\begin{minted}[fontsize=\small, breaklines, breakanywhere]{python}
from two_body.core.config import Config
from two_body.logic.ga import StreamingGA

ga = StreamingGA(Config())
population = ga.current_population()
print("Población inicial (primeros 3 individuos):", population[:3])
\end{minted}
\end{minipage}
\end{adjustbox}
\end{listing}

% ---------- fitness.py ----------
\subsection{\texttt{fitness.py}\-- Evaluación de Fitness}
\paragraph{Descripción general}
\texttt{FitnessEvaluator} calcula el fitness de cada candidato negando el exponente de Lyapunov reportado por \texttt{LyapunovEstimator}. Instrumenta cada evaluación con \texttt{time\_block}, utiliza un \texttt{HierarchicalCache} exacto (masas, horizonte, peso de periodicidad) y aplica penalizaciones por pérdida de cuasi-periodicidad cuando se solicita.

\subsubsection{\texttt{class FitnessEvaluator}}
\paragraph{Propósito} Proveer una función de costo escalable para el GA, reutilizando simulaciones cuando es posible, validando dominios y devolviendo metadatos (fitness, lambda, penalización, estado) cuando se solicitan detalles.

\subsubsection*{\texttt{evaluate\_batch~(candidates, horizon = ``short'') -> List[float]}}
\textbf{Propósito.} Evaluar un conjunto de candidatos devolviendo el fitness (mayor es mejor), reutilizando caché exacta, validando dominios y aplicando penalizaciones/telemetría según el horizonte solicitado.

\begin{table}[H]
  \centering
  \caption{Parómetros de \texttt{evaluate\_batch}}
  \small
  \begin{tabularx}{\textwidth}{@{}p{0.27\textwidth}p{0.33\textwidth}X@{}}
    \toprule
    \textbf{Parómetro} & \textbf{Tipo} & \textbf{Descripción} \\
    \midrule
    \texttt{candidates} & \texttt{List[Tuple[float, \dots]]} & Masas propuestos por el GA.\\
    \texttt{horizon} & \texttt{str} & Horizonte temporal (\texttt{``short''} o \texttt{``long''}).\\
    \texttt{epoch} & \texttt{int} & Identificador de época para telemetría; por defecto \texttt{-1}.\\
    \texttt{return\_details} & \texttt{bool} & Devuelve metadatos adicionales cuando es \texttt{True}.\\
    \bottomrule
  \end{tabularx}
\end{table}

\begin{table}[H]
  \centering
  \caption{Valor de retorno de \texttt{evaluate\_batch}}
  \small
  \begin{tabularx}{\textwidth}{@{}p{0.35\textwidth}X@{}}
    \toprule
    \textbf{Tipo} & \textbf{Descripción} \\
    \midrule
    \texttt{List[float] | Tuple[List[float], List[Dict[str, Any]]]} & Fitness por candidato (y opcionalmente detalles por individuo).\\
    \bottomrule
  \end{tabularx}
\end{table}

\begin{listing}[H]
\caption{Bloque de código de \texttt{evaluate\_batch}}
\begin{adjustbox}{max width=\textwidth}
\begin{minipage}{\linewidth}
{\setstretch{0.95}
\begin{adjustbox}{max width=\textwidth}
\begin{minipage}{\textwidth}
\begin{minted}[fontsize=\scriptsize, breaklines, breakanywhere]{python}
    def evaluate_batch(
        self,
        candidates: List[Tuple[float, ...]],
        horizon: str = "short",
        *,
        epoch: int = -1,
        return_details: bool = False,
    ) -> Union[List[float], Tuple[List[float], List[Dict[str, Any]]]]:
        import numpy as _np

        if len(candidates) == 0:
            return ([], []) if return_details else []

        t_end = self.cfg.t_end_short if horizon == "short" else self.cfg.t_end_long
        dt = self.cfg.dt
        sim_builder = ReboundSim(G=self.cfg.G, integrator=self.cfg.integrator)
        estimator = LyapunovEstimator()
        bounds = self.cfg.mass_bounds
        periodicity_weight = float(getattr(self.cfg, "periodicity_weight", 0.0) or 0.0)
        try:
            base_r0 = tuple(
                tuple(float(coord) for coord in vec) for vec in self.cfg.r0  # type: ignore[arg-type]
            )
            base_v0 = tuple(
                tuple(float(coord) for coord in vec) for vec in self.cfg.v0  # type: ignore[arg-type]
            )
        except Exception as exc:
            raise ValueError(
                "Config.r0 y Config.v0 deben contener vectores iterables numericos de longitud 3."
            ) from exc
        for name, base_vecs in (("r0", base_r0), ("v0", base_v0)):
            for idx, vec in enumerate(base_vecs):
                if len(vec) != 3:
                    raise ValueError(f"Config.{name}[{idx}] debe contener exactamente 3 componentes.")

        batch_id = self._next_batch_id()
        results: List[float] = []
        details: List[Dict[str, Any]] = []
        with time_block(
            "batch_eval",
            epoch=epoch,
            batch_id=batch_id,
            extra={"size": len(candidates), "horizon": horizon},
        ):
            for individual_id, masses_raw in enumerate(candidates):
                masses_tuple = tuple(float(m) for m in masses_raw)
                key_exact = (
                    tuple(round(m, 12) for m in masses_tuple)
                    + (horizon, round(periodicity_weight, 12))
                )
                timing_meta = {
                    "epoch": epoch,
                    "batch_id": batch_id,
                    "individual_id": individual_id,
                }
                with time_block(
                    "fitness_eval",
                    epoch=epoch,
                    batch_id=batch_id,
                    individual_id=individual_id,
                    extra={"horizon": horizon},
                ):
                    cached_payload = self.cache.get_exact(key_exact)
                    lam_val: Optional[float] = None
                    penalty_val: Optional[float] = None
                    status = "cache" if cached_payload is not None else "ok"
                    if cached_payload is not None:
                        if isinstance(cached_payload, dict):
                            fit = float(cached_payload.get("fitness", -float("inf")))
                            raw_lambda = cached_payload.get("lambda")
                            lam_val = float(raw_lambda) if raw_lambda is not None else None
                            raw_penalty = cached_payload.get("penalty")
                            penalty_val = float(raw_penalty) if raw_penalty is not None else None
                            status = str(cached_payload.get("status", status))
                        else:
                            fit = float(cached_payload)
                    else:
                        try:
                            masses = masses_tuple
                            if len(masses) != len(bounds):
                                raise ValueError(
                                    f"Cantidad de masas ({len(masses)}) no coincide con mass_bounds ({len(bounds)})."
                                )
                            for idx, m in enumerate(masses):
                                lo, hi = bounds[idx]
                                if m < lo or m > hi:
                                    raise ValueError(f"Masa[{idx}]={m} fuera de [{lo}, {hi}].")
                            if len(base_r0) < len(masses) or len(base_v0) < len(masses):
                                raise ValueError(
                                    "Config.r0 y Config.v0 deben definir al menos tantas entradas como masas."
                                )
                            r0 = base_r0[: len(masses)]
                            v0 = base_v0[: len(masses)]
                            sim = sim_builder.setup_simulation(masses=masses, r0=r0, v0=v0)
                            initial_r: Optional[_np.ndarray] = None
                            initial_v: Optional[_np.ndarray] = None
                            if periodicity_weight != 0.0:
                                initial_coords = []
                                initial_vels = []
                                for idx, particle in enumerate(sim.particles):
                                    if idx >= len(masses):
                                        break
                                    initial_coords.append(
                                        [float(particle.x), float(particle.y), float(particle.z)]
                                    )
                                    initial_vels.append(
                                        [float(particle.vx), float(particle.vy), float(particle.vz)]
                                    )
                                if initial_coords:
                                    initial_r = _np.asarray(initial_coords, dtype=float)
                                if initial_vels:
                                    initial_v = _np.asarray(initial_vels, dtype=float)
                            ctx = {
                                "sim": sim,
                                "dt": dt,
                                "t_end": t_end,
                                "masses": masses,
                                "timing_meta": timing_meta,
                            }
                            if len(masses) > 0:
                                ctx["m1"] = masses[0]
                            if len(masses) > 1:
                                ctx["m2"] = masses[1]
                            if len(masses) > 2:
                                ctx["m3"] = masses[2]
                            ret = estimator.mLCE(ctx, window=t_end, timing_meta=timing_meta)
                            lam = float(ret.get("lambda", _np.inf))
                            lam_val = lam
                            if not _np.isfinite(lam):
                                status = "nonfinite_lambda"
                                self.logger.debug(
                                    "Lyapunov no finito para %s (horizon=%s): %s",
                                    masses_tuple,
                                    horizon,
                                    lam,
                                )
                                fit = -float("inf")
                            else:
                                penalty = 0.0
                                if (
                                    periodicity_weight != 0.0
                                    and initial_r is not None
                                    and initial_v is not None
                                ):
                                    try:
                                        rf_list: list[list[float]] = []
                                        vf_list: list[list[float]] = []
                                        for idx, particle in enumerate(sim.particles):
                                            if idx >= len(masses):
                                                break
                                            rf_list.append(
                                                [float(particle.x), float(particle.y), float(particle.z)]
                                            )
                                            vf_list.append(
                                                [float(particle.vx), float(particle.vy), float(particle.vz)]
                                            )
                                        rf_arr = _np.asarray(rf_list, dtype=float)
                                        vf_arr = _np.asarray(vf_list, dtype=float)
                                        if rf_arr.shape == initial_r.shape and vf_arr.shape == initial_v.shape:
                                            delta_r = _np.linalg.norm(rf_arr - initial_r, axis=1).sum()
                                            delta_v = _np.linalg.norm(vf_arr - initial_v, axis=1).sum()
                                            penalty = float(delta_r + delta_v)
                                        if not _np.isfinite(penalty):
                                            self.logger.debug(
                                                "Penalizacion no finita: delta_r=%.3e, delta_v=%.3e",
                                                delta_r,
                                                delta_v,
                                            )
                                            penalty = float("inf")
                                        else:
                                            self.logger.debug(
                                                "delta_r=%.6f, delta_v=%.6f, penalty=%.6f",
                                                delta_r,
                                                delta_v,
                                                penalty,
                                            )
                                            penalty = min(penalty, 1e6)
                                    except Exception:
                                        penalty = 0.0
                                penalty_val = penalty
                                fit = -lam - periodicity_weight * penalty
                        except Exception as exc:
                            status = "error"
                            self.logger.debug(
                                "Fallo evaluando candidato %s (horizon=%s): %s",
                                masses_tuple,
                                horizon,
                                exc,
                                exc_info=True,
                            )
                            fit = -float("inf")
                            lam_val = lam_val if lam_val is not None else None
                        cache_payload = {
                            "fitness": float(fit),
                            "lambda": float(lam_val) if lam_val is not None else None,
                            "penalty": float(penalty_val) if penalty_val is not None else None,
                            "status": status,
                        }
                        self.cache.set_exact(key_exact, cache_payload)

                results.append(float(fit))
                details.append(
                    {
                        "masses": masses_tuple,
                        "fitness": float(fit),
                        "lambda": float(lam_val) if lam_val is not None else None,
                        "penalty": float(penalty_val) if penalty_val is not None else None,
                        "cached": cached_payload is not None,
                        "status": status,
                        "horizon": horizon,
                    }
                )

        if return_details:
            return results, details
        return results
\end{minted}
\end{minipage}
\end{adjustbox}


}
\end{minipage}
\end{adjustbox}
\end{listing}



\begin{listing}[H]
\caption{Ejemplo de evaluación en lote}
\begin{adjustbox}{max width=\textwidth}
\begin{minipage}{\textwidth}
\begin{minted}[fontsize=\small, breaklines, breakanywhere]{python}
from two_body.core.cache import HierarchicalCache
from two_body.core.config import Config
from two_body.logic.fitness import FitnessEvaluator

cfg = Config()
cache = HierarchicalCache()
evaluator = FitnessEvaluator(cache, cfg)
scores, details = evaluator.evaluate_batch(
    [(1.0, 1.0)],
    horizon="short",
    epoch=0,
    return_details=True,
)
print("Fitness:", scores, "Detalles:", details)
\end{minted}
\end{minipage}
\end{adjustbox}
\end{listing}

% ---------- parameters.py ----------
\subsection{\texttt{parameters.py} \-- Modificación de Parómetros}
\paragraph{Descripción general}
\texttt{ParameterModifier} encapsula la estrategia de re-seeding cuando el GA se estanca. Centraliza la lógica que antes residóa en el controlador monolítico y permite ajustar la exploración local (radio).

\subsubsection{\texttt{class ParameterModifier}}
\paragraph{Propósito} Ajustar el radio de exploración y reposicionar al GA alrededor del mejor candidato conocido tras periodos de estancamiento.

\subsubsection*{\texttt{on\_stagnation~(ga, best, radius) -> float}}
\textbf{Propósito.} Reseedear el GA alrededor del mejor candidato y devolver el nuevo radio de exploración escalado por \texttt{radius\_decay}.

\begin{table}[H]
  \centering
  \caption{Parómetros de \texttt{on\_stagnation}}
  \small
  \begin{tabularx}{\textwidth}{@{}p{0.27\textwidth}p{0.33\textwidth}X@{}}
    \toprule
    \textbf{Parómetro} & \textbf{Tipo} & \textbf{Descripción} \\
    \midrule
    \texttt{ga} & \texttt{StreamingGA} & Instancia del GA a la que se le aplicaró el reseed. \\
    \texttt{best} & \texttt{Tuple[float, \dots] | None} & Mejor candidato conocido; if es \texttt{None}, no se aplica reseed. \\
    \texttt{radius} & \texttt{float} & Radio actual de exploración local. \\
    \bottomrule
  \end{tabularx}
\end{table}

\begin{table}[H]
  \centering
  \caption{Valor de retorno de \texttt{on\_stagnation}}
  \small
  \begin{tabularx}{\textwidth}{@{}p{0.35\textwidth}X@{}}
    \toprule
    \textbf{Tipo} & \textbf{Descripción} \\
    \midrule
    \texttt{float} & Nuevo radio recomendado tras aplicar \texttt{cfg.radius\_decay}. \\
    \bottomrule
  \end{tabularx}
\end{table}

\begin{listing}[H]
\caption{Bloque de código de \texttt{on\_stagnation}}
\begin{adjustbox}{max width=\textwidth}
\begin{minipage}{\textwidth}
\begin{minted}[fontsize=\small, breaklines, breakanywhere]{python}
class ParameterModifier:
    def on_stagnation(
        self,
        ga: StreamingGA,
        best: Optional[Tuple[float, ...]],
        radius: float,
    ) -> float:
        if best is None:
            return radius
        ga.reseed_around(best, radius)
        return radius * self.cfg.radius_decay
\end{minted}
\end{minipage}
\end{adjustbox}
\end{listing}

\begin{listing}[H]
\caption{Ejemplo de reseed tras estancamiento}
\begin{adjustbox}{max width=\textwidth}
\begin{minipage}{\textwidth}
\begin{minted}[fontsize=\small, breaklines, breakanywhere]{python}
from two_body.core.config import Config
from two_body.logic.ga import StreamingGA
from two_body.logic.parameters import ParameterModifier

cfg = Config()
ga = StreamingGA(cfg)
modifier = ParameterModifier(cfg)
new_radius = modifier.on_stagnation(ga, (1.0, 1.0), cfg.local_radius)
print("Nuevo radio:", new_radius)
\end{minted}
\end{minipage}
\end{adjustbox}
\end{listing}

% ---------- __init__.py ----------
\subsection{\texttt{\_\_init\_\_.py} \-- API del Módulo}
\paragraph{Descripción general}
El archivo \texttt{\_\_init\_\_.py} re-exporta las clases principales para facilitar la importación desde \texttt{two\_body.logic}. Asó, los consumidores pueden escribir \texttt{from two\_body.logic import StreamingGA} sin exponer detalles internos.

\begin{listing}[H]
\caption{Bloque de código de \texttt{logic.\_\_init\_\_.py}}
\begin{adjustbox}{max width=\textwidth}
\begin{minipage}{\textwidth}
\begin{minted}[fontsize=\small, breaklines, breakanywhere]{python}
"""
Capa logica: motor de optimizacion y adaptacion de parametros.
"""

from .fitness import FitnessEvaluator  # noqa: F401
from .ga import StreamingGA  # noqa: F401
from .parameters import ParameterModifier  # noqa: F401
from .controller import ContinuousOptimizationController  # noqa: F401

__all__ = ["FitnessEvaluator", "StreamingGA", "ParameterModifier", "ContinuousOptimizationController"]
\end{minted}
\end{minipage}
\end{adjustbox}
\end{listing}

\begin{listing}[H]
\caption{Ejemplo de importación desde la API póblica}
\begin{adjustbox}{max width=\textwidth}
\begin{minipage}{\textwidth}
\begin{minted}[fontsize=\small, breaklines, breakanywhere]{python}
from two_body.logic import ContinuousOptimizationController, StreamingGA
\end{minted}
\end{minipage}
\end{adjustbox}
\end{listing}

% --- Fin de capítulo ---
